\chapter{项目六：命令行文件索引器}

\section{项目目标}

文件索引器扫描一个目录，统计每种扩展名的文件数量和总字节数，输出排序后的报表。这个项目覆盖文件系统遍历、错误处理、路径处理、容器聚合、排序、命令行参数和性能边界。

目标命令：

\begin{lstlisting}[language=bash,caption={文件索引器用法}]
fileindex --root ./src --top 20
\end{lstlisting}

输出示例：

\begin{lstlisting}[caption={文件索引器输出}]
.cpp 128 files 240000 bytes
.hpp 96 files 120000 bytes
.md 12 files 34000 bytes
\end{lstlisting}

\section{数据模型}

\begin{lstlisting}[language=C++,caption={扩展名统计}]
struct ExtensionStats {
    std::uint64_t file_count = 0;
    std::uint64_t total_bytes = 0;
};

using ExtensionIndex = std::unordered_map<std::string, ExtensionStats>;
\end{lstlisting}

这里用 \code{std::uint64_t}，因为文件数量和字节数可能超过 \code{int} 范围。类型选择来自值域，不是随手写。

\section{路径和扩展名规则}

扩展名规则要先写清楚：

\begin{itemize}
  \item 没有扩展名的文件归为 \code{<none>}。
  \item 扩展名保留点号，例如 \code{.cpp}。
  \item 大小写是否合并由项目决定；本项目统一转小写。
\end{itemize}

实现：

\begin{lstlisting}[language=C++,caption={提取扩展名}]
std::string normalized_extension(const std::filesystem::path& path)
{
    std::string extension = path.extension().string();
    if (extension.empty()) {
        return "<none>";
    }
    for (char& ch : extension) {
        ch = static_cast<char>(std::tolower(static_cast<unsigned char>(ch)));
    }
    return extension;
}
\end{lstlisting}

这个函数可以单独测试，不需要扫描真实目录。

\section{遍历目录}

\begin{lstlisting}[language=C++,caption={扫描目录}]
ExtensionIndex build_index(const std::filesystem::path& root)
{
    ExtensionIndex index;
    std::error_code error;
    std::filesystem::recursive_directory_iterator it{
        root,
        std::filesystem::directory_options::skip_permission_denied,
        error};
    std::filesystem::recursive_directory_iterator end;

    for (; it != end; it.increment(error)) {
        if (error) {
            error.clear();
            continue;
        }
        if (!it->is_regular_file(error) || error) {
            error.clear();
            continue;
        }

        const std::uint64_t size = it->file_size(error);
        if (error) {
            error.clear();
            continue;
        }
        ExtensionStats& stats = index[normalized_extension(it->path())];
        ++stats.file_count;
        stats.total_bytes += size;
    }
    return index;
}
\end{lstlisting}

这里用 \code{std::error_code} 避免权限问题直接抛异常终止扫描。项目选择是：跳过无法访问的文件。生产工具应记录警告；教材版本先保留核心流程。

\section{排序报表}

\begin{lstlisting}[language=C++,caption={排序扩展名统计}]
struct ExtensionRow {
    std::string extension;
    ExtensionStats stats;
};

std::vector<ExtensionRow> sort_index(const ExtensionIndex& index)
{
    std::vector<ExtensionRow> rows;
    rows.reserve(index.size());
    for (const auto& [extension, stats] : index) {
        rows.push_back(ExtensionRow{extension, stats});
    }

    std::ranges::sort(rows, [](const ExtensionRow& lhs, const ExtensionRow& rhs) {
        if (lhs.stats.total_bytes != rhs.stats.total_bytes) {
            return lhs.stats.total_bytes > rhs.stats.total_bytes;
        }
        return lhs.extension < rhs.extension;
    });
    return rows;
}
\end{lstlisting}

排序规则必须稳定可解释：总字节数大的排前面，字节数相同按扩展名字典序。这样测试输出不会因为哈希表顺序变化而随机变化。

\section{输出层}

\begin{lstlisting}[language=C++,caption={写报表}]
void write_report(std::ostream& output, std::span<const ExtensionRow> rows, std::size_t limit)
{
    const std::size_t count = std::min(rows.size(), limit);
    for (std::size_t i = 0; i < count; ++i) {
        const ExtensionRow& row = rows[i];
        output << row.extension << " "
               << row.stats.file_count << " files "
               << row.stats.total_bytes << " bytes\n";
    }
}
\end{lstlisting}

输出函数接受 \code{std::ostream&}，因此可以用字符串流测试，不必真正打印到终端。

\section{性能边界}

目录遍历是 I/O 密集型。优化前先明确瓶颈：

\begin{itemize}
  \item 文件数量很大时，遍历系统调用是主要成本。
  \item 扩展名种类通常不多，哈希表不是瓶颈。
  \item 输出 top k 可以全排序，也可以用堆；当扩展名种类少时全排序足够。
  \item 并行遍历可能受磁盘和文件系统限制，不一定更快。
\end{itemize}

第一版不要急着并行。先让错误处理和结果稳定。等测量证明遍历慢，再考虑多线程扫描。

\section{测试计划}

测试可以创建临时目录：

\begin{itemize}
  \item 空目录输出为空。
  \item 没有扩展名的文件归为 \code{<none>}。
  \item \code{.CPP} 和 \code{.cpp} 合并。
  \item 多个扩展名按总字节数排序。
  \item 权限错误不让整个扫描崩溃。
\end{itemize}

临时目录要用 RAII 管理，测试结束自动删除。这样测试不会污染工作区。

\begin{exercisebox}
给文件索引器增加 \code{--min-size N}，只统计至少 \code{N} 字节的文件。要求参数解析拒绝负数和非数字；扫描逻辑不关心命令行字符串，只接收已经解析好的整数。
\end{exercisebox}

\section{扫描策略要先定义错误政策}

文件系统遍历会遇到权限不足、符号链接循环、文件在扫描中被删除、路径太长等问题。如果不先定义错误政策，代码就会在异常和忽略之间摇摆。本项目采用：

\begin{itemize}
  \item 根目录打不开：命令失败。
  \item 子目录权限不足：跳过并记录警告。
  \item 单个文件读取大小失败：跳过并记录警告。
  \item 符号链接：第一版不跟随。
\end{itemize}

这些政策会影响接口。只返回 \code{ExtensionIndex} 不够，还应返回警告：

\begin{lstlisting}[language=C++,caption={索引结果带警告}]
struct IndexWarning {
    std::filesystem::path path;
    std::string message;
};

struct IndexResult {
    ExtensionIndex index;
    std::vector<IndexWarning> warnings;
};
\end{lstlisting}

这样 CLI 可以把警告写到错误流，测试也能断言权限问题没有让整个扫描失败。

\section{遍历循环的证据字段}

\code{recursive_directory_iterator} 使用 \code{error_code} 时，每一步都要检查。可以把循环理解成三段：

\begin{enumerate}
  \item 迭代器推进是否成功。
  \item 当前条目是否是普通文件。
  \item 文件大小是否可读。
\end{enumerate}

每段失败都产生不同警告。不要把所有失败都简单 \code{continue}，否则用户不知道为什么统计少了文件。教材里可以先简化，但工程版本必须留下错误证据。

\section{top k 是否需要堆}

输出前 \code{k} 个扩展名时有两种做法。第一种是全部排序，复杂度 \complexity{O(m \log m)}，其中 \code{m} 是扩展名种类数。第二种是维护大小为 \code{k} 的堆，复杂度 \complexity{O(m \log k)}。听起来堆更高级，但文件扩展名种类通常很少，全部排序简单、稳定、可读，完全足够。

这个例子训练算法选择：不要因为知道更复杂的数据结构就立刻使用。先估计 \code{m} 的规模。如果 \code{m} 只有几十或几百，排序成本远小于文件系统遍历成本。真正热点在系统调用和磁盘 I/O。

\section{min-size 参数贯穿三层}

增加 \code{--min-size N} 时，要分三层：

\begin{longtable}{p{0.22\linewidth}p{0.32\linewidth}p{0.34\linewidth}}
\toprule
层 & 输入 & 责任 \\
\midrule
参数解析 & 字符串 \code{N} & 完整解析非负整数 \\
扫描配置 & \code{std::uint64_t min_size} & 保存已经验证的阈值 \\
扫描逻辑 & 文件大小 & 小于阈值则跳过 \\
\bottomrule
\end{longtable}

扫描函数不应该知道 \code{"--min-size"} 这个字符串。它只接收配置对象：

\begin{lstlisting}[language=C++,caption={扫描配置对象}]
struct IndexOptions {
    std::filesystem::path root;
    std::uint64_t min_size = 0;
    std::size_t top = 20;
};
\end{lstlisting}

这样命令行、配置文件或图形界面都可以复用同一个扫描逻辑。

\section{临时目录测试}

文件索引器需要真实文件系统测试，但要控制范围。测试可以创建临时目录：

\begin{lstlisting}[language=C++,caption={临时目录测试思路}]
TemporaryDirectory temp;
write_file(temp.path() / "a.cpp", "1234");
write_file(temp.path() / "b.HPP", "12");
write_file(temp.path() / "README", "x");

const IndexResult result = build_index(IndexOptions{.root = temp.path()});
assert(result.index.at(".cpp").file_count == 1);
assert(result.index.at(".hpp").file_count == 1);
assert(result.index.at("<none>").file_count == 1);
\end{lstlisting}

\code{TemporaryDirectory} 应使用 RAII：构造时创建目录，析构时删除。这样测试失败时也尽量清理现场。不要把测试文件散落到仓库根目录。

\section{输出稳定性}

因为 \code{unordered_map} 遍历顺序不稳定，任何面向用户或测试的输出都必须先排序。排序规则要覆盖平局情况：字节数相同按扩展名升序，扩展名相同不可能出现。只要平局规则不完整，测试就可能在不同平台或不同运行中随机失败。

\section{项目验收清单}

文件索引器完成时至少验收：

\begin{itemize}
  \item 空目录输出为空。
  \item 大小写扩展名合并。
  \item 无扩展名归入固定名称。
  \item top 限制生效。
  \item min-size 拒绝非法参数。
  \item 单文件错误产生警告而不是崩溃。
  \item 输出顺序稳定。
\end{itemize}

这些验收项覆盖了文件系统项目最容易出错的边界：路径、错误、排序、参数和可重复测试。
